\chapter{Introduction}
\label{chap:intro}

\section{Context and Motivation}

SQLite is the world's most widely deployed database engine. Unlike client-server database systems, SQLite operates as a self-contained, serverless library embedded directly into the host application. This design makes it the default database backend for every major smartphone operating system, every mainstream web browser, the Python standard library, and an extensive range of embedded and IoT devices. With an estimated three trillion active installations at any given time \cite{sqlite}, a single exploitable vulnerability in SQLite threatens not only the applications that use it directly but also the entire ecosystem of downstream software that depends on it. The attack surface is accordingly vast: any device, application, or service that processes untrusted SQL — whether from network input, user-supplied queries, or parsed file formats — is a potential exploitation vector.

The severity of this exposure is not merely theoretical. Between 2019 and 2020, six CVEs were filed against SQLite versions 3.30.1 through 3.32.2. CVE-2019-19646 \cite{cve201919646} exposes an infinite loop in aggregate query processing. CVE-2020-13434 \cite{cve202013434} triggers an integer overflow inside the built-in \texttt{printf} implementation. CVE-2020-9327 \cite{cve20209327} involves an uninitialized pointer in generated column handling. CVE-2020-13435 \cite{cve202013435} causes a null pointer dereference. CVE-2020-13871 \cite{cve202013871} is a use-after-free in window function processing. CVE-2020-15358 \cite{cve202015358} produces a heap buffer overread. The fact that such defects persisted across multiple release cycles of a mature, heavily tested, and widely audited codebase underscores a fundamental limitation of conventional testing methodologies: manual code review and unit testing do not scale to the combinatorial space of SQL statement sequences that exercise deep query optimizer, window function, and virtual table logic. Systematic automated approaches are necessary.

Fuzzing is one of the most productive techniques available for automated vulnerability detection \cite{fuzzingsurvey}. Coverage-guided greybox fuzzers such as AFL \cite{afl} and libFuzzer \cite{libfuzzer} instrument the target binary to collect branch coverage feedback, then use that feedback to steer mutation toward unexplored program regions. This feedback loop has proven highly effective for file parsers, network protocol implementations, and other software that consumes loosely structured inputs. However, byte-level mutation fuzzers encounter a fundamental barrier when applied to programs that expect highly structured inputs such as SQL: random bit flips and splice mutations produce inputs that are rejected immediately by the parser, never exercising the complex semantic layers — query planning, expression evaluation, window frame computation — where the most interesting vulnerabilities reside. The parser acts as a filter that prevents coverage-guided mutation from reaching the deep logic that matters.

Grammar-based fuzzing addresses this structural barrier by generating inputs from a context-free grammar, guaranteeing syntactic validity before the input reaches the target program. Nautilus \cite{nautilus} was the first system to combine grammar-based generation with AFL-style coverage feedback, achieving both structural validity and coverage guidance. Superion \cite{superion} extends AFL with grammar-aware tree mutations for JavaScript and XML targets. Grimoire \cite{grimoire} infers structural patterns from successful inputs without requiring an explicit grammar specification. Within the DBMS testing domain specifically, Squirrel \cite{squirrel} adds semantic validity awareness to SQL mutation, SQLsmith \cite{sqlsmith} generates random valid SQL queries without coverage feedback, SQLRight \cite{sqlright} leverages side-channel analysis to detect logic bugs, and Griffin \cite{griffin} performs grammar-free DBMS fuzzing through input reconstruction. Each of these systems advances the state of the art in a specific dimension: semantic validity, structural inference, or coverage efficiency.

A limitation shared by all these approaches is that the grammar itself receives little systematic engineering attention. Existing grammar-based fuzzers for SQL typically use general-purpose SQL grammars that treat all constructs with equal weight, making no distinction between a trivial \texttt{SELECT 1} and a window function with a complex frame specification nested inside a correlated subquery. A SELECT statement with a simple WHERE clause competes on equal footing with the precise structural patterns — generated columns with expression constraints, recursive common table expressions, FTS auxiliary function calls — that are known to exercise the code paths where CVEs were found. No existing system encodes domain knowledge about which SQL structural patterns are most likely to trigger specific vulnerability classes; no existing system reasons about grammar design as a first-class engineering problem. This gap motivates a focused investigation of grammar engineering as the primary lever for improving vulnerability discovery effectiveness in grammar-based DBMS fuzzing.


\section{Problem Statement}

This thesis investigates whether a carefully engineered grammar based on structural primitives — production rules that encode SQL patterns associated with known vulnerability classes — can improve the effectiveness of grammar-based fuzzing for DBMS vulnerability discovery. The structural primitives approach decomposes SQL into layered non-terminals that encode vulnerability-triggering patterns without embedding proof-of-concept inputs directly in grammar rules, thereby preserving the compositional freedom that allows the fuzzer to discover novel variants.

This investigation is guided by two research questions:

\begin{itemize}
    \item \textbf{RQ1:} Can DBMS-Nautilus rediscover existing CVEs in older SQLite versions through compositional grammar generation, without embedding proof-of-concept inputs in the grammar?
    \item \textbf{RQ2:} Can DBMS-Nautilus discover previously unreported bugs in SQLite beyond the known CVEs?
\end{itemize}

RQ1 establishes whether the structural primitives methodology is sufficient to reproduce known ground-truth vulnerabilities through emergent composition — the fuzzer combining grammar rules in ways the grammar designer did not explicitly specify. RQ2 probes whether the same methodology generalizes beyond the known vulnerability landscape to surface defects that have not yet been publicly reported.


\section{Objectives and Scope}

This thesis pursues three concrete objectives. The first is to design a structural primitives grammar methodology for SQL fuzzing. Rather than enumerating SQL constructs uniformly, this methodology decomposes SQL into layered non-terminal hierarchies — schema-setup operators, stress-query patterns, validation operations, boundary function calls — where each production rule encodes a structural pattern associated with a vulnerability-triggering code path. The critical constraint is that no rule may contain a complete proof-of-concept input: the grammar defines structural affordances, and the fuzzer's compositional generation is responsible for assembling those affordances into crash-inducing inputs.

The second objective is to implement DBMS-Nautilus as a complete, end-to-end fuzzing system. This involves extending the Nautilus fuzzer with the structural grammar, implementing an AFL-compatible fork server harness instrumented with ASan and UBSan for four target SQLite versions, and building a triage pipeline that automates crash deduplication, severity classification, and CVE correlation.

The third objective is to conduct an empirical evaluation across the four SQLite versions (3.30.1, 3.31.1, 3.32.0, and 3.32.2) that contain the six confirmed CVEs, measuring both CVE rediscovery rate and the discovery of previously unreported bugs.

The scope of this work is bounded in several important ways. The grammar evolves through four versions (v3.0 with 449 rules through v3.3 with 520 rules), with each version targeting a specific deficiency identified by coverage analysis or crash evidence. The evaluation focuses on a single DBMS, SQLite, chosen for its ubiquity and its well-documented CVE history. The grammar is a context-free grammar only; no semantic constraints such as type consistency or schema-aware expression generation are enforced.


\section{Contributions}

This thesis makes three primary contributions to the field of grammar-based DBMS fuzzing.

The first contribution is the structural primitives grammar methodology with zero proof-of-concept contamination. The methodology defines a principled approach to grammar design in which production rules encode vulnerability-class-relevant structural patterns — \texttt{Schema-Setup}, \texttt{Stress-Query}, \texttt{Validation-Op}, \texttt{Boundary-Func-Call} and related non-terminals — without any rule containing a complete exploit. This design choice is deliberate: if a grammar rule directly encodes a known CVE proof-of-concept, the fuzzer's rediscovery of that CVE is a retrieval task, not a discovery task. The structural primitives methodology preserves compositional freedom, allowing the fuzzer to rediscover known CVEs through emergent combination and to discover novel bugs not anticipated by the grammar designer.

The second contribution is an iterative grammar evolution methodology driven by empirical evidence. Starting from a v3.0 baseline of 449 rules, grammar evolution proceeds through v3.1 (addressing FTS rule weight dominance that suppressed other coverage), v3.2 (adding structural patterns for CVEs not yet reached), and v3.3 (expanding JSON and JSONB operator coverage to close identified gaps). Each evolution step is justified by coverage profiling data and crash evidence, establishing a feedback loop between empirical measurement and grammar design.

The third contribution is an empirical demonstration of the methodology's effectiveness across four SQLite versions containing six confirmed CVEs. DBMS-Nautilus rediscovers three of the six target CVEs through compositional grammar generation. Beyond CVE rediscovery, the system discovers seven previously unreported bug classes across 89 unique crash instances spanning all four target versions, including memory corruption, null pointer dereferences, and signed integer overflow behaviors identified through UBSan instrumentation.


\section{Thesis Organization}

The remainder of this thesis is organized as follows. Chapter~\ref{chap:background} provides background on coverage-guided and grammar-based fuzzing, reviews related DBMS testing tools, describes the sanitizer infrastructure (ASan and UBSan) used as crash oracles, and gives a brief overview of the SQLite architecture and the specific code regions implicated in the target CVEs. Chapter~\ref{chap:method} presents the DBMS-Nautilus system design in full: the overall architecture, the structural primitives grammar methodology and its non-terminal hierarchy, the iterative grammar evolution from v3.0 to v3.3, the AFL-compatible fork server harness, and the automated triage pipeline. Chapter~\ref{chap:experiments} reports the empirical evaluation, structured around the two research questions: Section~4.1 addresses RQ1 through CVE rediscovery experiments with per-CVE analysis, and Section~4.2 addresses RQ2 through new bug discovery results with crash classification and severity assessment; a coverage analysis section connects grammar design decisions to observed coverage trajectories. The final chapter synthesizes the findings, discusses limitations of the structural primitives approach, and identifies directions for future work including semantic constraint integration, multi-DBMS generalization, and automated grammar evolution.
